\chapter{Black holes and holography}\label{ch:bh}

A condensed treatment of the same ideas for black holes: large $N$ and emergent type III$_1$, the black-hole crossed product, entanglement wedges and error correction, and Hawking radiation in curved spacetime.

\section{Large $N$ and emergent type III$_1$}\label{ch:large_N}
\skippable

\begin{physmotiv}
In AdS/CFT, a boundary theory at large $N$ describes bulk gravity. The boundary algebra of a single-sided region is type I in finite $N$ regularizations but, in the large-$N$ limit around a thermal state (the eternal black hole), the algebra of single-trace operators becomes a type III$_1$ algebra: an emergent type III$_1$ factor with an emergent time, the Schwarzschild time of the exterior. This is the black hole counterpart of the type III$_1$ static patch algebra, and the starting point for the crossed-product treatment of the black hole. Papers: Leutheusser and Liu on emergent times and subalgebra--subregion duality; Witten, and Chandrasekaran--Penington--Witten on large $N$ algebras.
\end{physmotiv}

\subsection{The two-sided black hole and the thermofield double}

The eternal AdS black hole is dual to the thermofield double of two copies of the CFT, $\ket{\rm TFD}=Z^{-1/2}\sum\ee^{-\beta E_n/2}\ket n_L\ket n_R$ (\cref{ch:tomita}). At finite $N$ the right CFT algebra is $\Bnd{\HH_R}$, type I, and the TFD is a purification of the thermal state with $\Delta=\ee^{-\beta(H_R-H_L)}$. In the strict large-$N$ limit, define the algebra $\mathcal M_R$ generated by (bounded functions of) single-trace operators smeared in time, in the GNS representation of the thermal state. Because the thermal state has an infinite number of low-lying single-particle excitations of the bulk fields in the $N\to\infty$ limit, the modular spectrum becomes continuous, and the algebra is expected to be type III$_1$ (Leutheusser--Liu \cite{LeutheusserLiu2022}), the bulk free fields in the exterior being a generalized free field theory at $N=\infty$ with a continuous spectrum for $\Delta$, with the boundary time translation as outer modular flow. In this limit, boundary time $\leftrightarrow$ modular time up to $\beta$, as in \cref{ch:rs_bw}.

\subsection{What this tells us}
\begin{itemize}
\item \textbf{Emergent horizon.} The type III$_1$ factor encodes the bulk exterior; its commutant is the bulk exterior of the other side, the algebra of $L$ (Haag duality at large $N$).
\item \textbf{Why type III.} The bulk exterior has horizon-ward modes at all energy scales. Type III$_1$ is the expected type of the exterior bulk algebra for the same reason as in \cref{ch:typeIII_local}.
\item \textbf{Subalgebra--subregion duality.} Different boundary subalgebras correspond to different bulk regions; this is the algebraic form of entanglement-wedge reconstruction (\cref{ch:ew_qec}).
\end{itemize}

\begin{whatwrong}
Type III$_1$ appears only in the large-$N$ limit: at finite $N$ the boundary algebra is type I. The statement that this is type III$_1$ relies on properties of generalized free field theory at $N=\infty$ and is argued, not proved for any interacting CFT. The entropy of the boundary theory at finite $N$ is finite; the large-$N$ type III$_1$ limit loses it in the same way as the continuum limit of a lattice.
\end{whatwrong}

\subsection*{Exercises}
\begin{exercise}\label{ex:large_N:1}
For a single bulk free scalar in the BTZ-like exterior with quasinormal-mode density, argue that the thermal state of the single-trace algebra has $\Delta$ with spectrum $\ee^{-\beta\omega}$, $\omega\in\R$, hence continuous.
\end{exercise}
\begin{exercise}\label{ex:large_N:2}
Explain why the commutant of the right single-trace algebra in the TFD GNS space is the left single-trace algebra.
\end{exercise}

\section{Crossed product for black holes}\label{ch:bh_crossed}
\skippable

\begin{physmotiv}
For the large-$N$ black hole, the boundary Hamiltonian $H$ (the ADM energy) is a physical operator, not a constraint. Witten's observation is that this is precisely what turns the type III$_1$ bulk exterior algebra into a type II$_\infty$ one: the algebra of the boundary at large $N$ is generated by the bulk exterior fields \emph{and} $H$ (dressing the bulk operators), i.e.\ it is a crossed product. Chandrasekaran, Penington and Witten computed the entropy and showed it is the generalized entropy.
\end{physmotiv}

\subsection{Witten's construction}

Let $\A$ be the type III$_1$ algebra of bulk fields in the exterior of the eternal black hole in the Hartle--Hawking state, with modular flow $\Ad\ee^{-\ii\beta tH_{\rm bulk}}$-type (the Schwarzschild boost). Gauge invariance under diffeomorphisms that act at the boundary means that bulk operators must be dressed to the boundary, and the boundary Hamiltonian $H$ (which has a fluctuating spectrum in the large-$N$ limit of order $N^0$ around the classical mass) is a new operator in the algebra. Writing $H=\frac{1}{\beta}K+\ldots$ with $K$ the modular Hamiltonian of the bulk fields up to a constant, the algebra is
\begin{equation}
\A\rtimes_\sigma\R\;\cong\;\{\hat a,\ H\ \text{(shifted)}\}'',
\end{equation}
with the physical Hamiltonian playing the role of $X=K+x$ in \cref{ch:crossed}. It is type II$_\infty$ (the energy spectrum is unbounded above but the trace is $\ee^{\beta H}$-weighted, \cref{ch:takesaki}). \emph{There is no maximum-entropy state}\index{maximum-entropy state}: the black hole is not in equilibrium with a heat bath in the same sense as de Sitter with a bounded-below observer; instead the entropy of a state is defined up to an additive constant and is unbounded above, as expected for a black hole that can grow.

\subsection{Entropy}

Chandrasekaran, Penington and Witten \cite{CPW2022} showed that the entropy of semiclassical states of the crossed product agrees, up to a state-independent constant, with the generalized entropy $\Sgen=A/4G_N+S_{\rm bulk}$; this is the first and cleanest derivation of the generalized entropy from a von Neumann entropy of a type II algebra (the formula derived in a finite-dimensional model in \cref{ch:entropy_II}).

\begin{center}
\renewcommand{\arraystretch}{1.3}
\resizebox{\ifdim\width>\linewidth\linewidth\else\width\fi}{!}{%
\begin{tabular}{@{}lll@{}}
\toprule
 & Black hole (large $N$) & de Sitter (observer)\\
\midrule
Clock & boundary Hamiltonian $H$, spectrum unbounded above & observer energy $q\ge0$\\
Algebra & type II$_\infty$ & type II$_1$\\
Max entropy state & none & empty dS\\
Trace weight & $\ee^{\beta H}$ & $\ee^{-\beta q}$\\
\bottomrule
\end{tabular}}
\end{center}

\subsection{Relation to the de Sitter construction}
CLPW devote their \S5 to revisiting the type II$_\infty$ algebra of the black hole of \cite{CPW2022}, with an alternative derivation that makes the analogy with the Hilbert space of two observers in de Sitter space explicit. The structural dictionary is: enlarge $\HH\to\HH\otimes L^2(\R)$ and add one generator $H+x$; this is the de Sitter construction with $q=-x$, \emph{except} that de Sitter carries the constraint $q\ge0$ and the black hole has no analogue. Without the projection the trace weight $\ee^{x}\dd x$ has infinite total mass, so the algebra is II$_\infty$ and there is no maximum-entropy state. Near $N=\infty$ the exterior algebra of the black hole is type II$_\infty$ \cite{CPW2022}; CLPW note that the black hole entropy can grow indefinitely with the mass, which is why the missing bound is physically sensible.

\begin{whatwrong}
The black hole construction has a physical $H$ already in the boundary theory. In de Sitter there is none; the observer must be added. The algebraic structure is the same; the physical origin of the clock is different, and so is the sign structure that decides II$_\infty$ versus II$_1$.
\end{whatwrong}

\subsection*{Exercises}
\begin{exercise}\label{ex:bh_crossed:1}
Using \cref{ch:takesaki}, show that the trace on the black hole algebra is $\tau(\hat a)=\int\dd x\,\ee^{x}\ldots$ with $x\sim\beta(H-\bar H)$, and that $\tau(\id)=\infty$ because $x$ has no upper bound.
\end{exercise}
\begin{exercise}\label{ex:bh_crossed:2}
Compare the two-column table above: which sign flip converts $\ee^{\beta H}$ to $\ee^{-\beta q}$ in the trace weight, and what does this say about the direction in which energy can be added?
\end{exercise}

\section{Entanglement wedges and error correction}\label{ch:ew_qec}
\skippable

\begin{physmotiv}
The algebraic language also clarifies holographic subregion duality. A boundary subregion's algebra, in the large-$N$ code subspace, equals the algebra of its entanglement wedge: this is operator-algebra quantum error correction, where bulk operators are logical operators and the boundary algebra of a region is a ``code'' for them. The crossed-product entropy then becomes the algebraic origin of the quantum extremal surface prescription.
\end{physmotiv}

\subsection{Subregion duality as an algebra statement}

Jafferis, Lewkowycz, Maldacena and Suh (JLMS \cite{JLMS2016}) showed that, to leading order in $1/N$, the boundary modular Hamiltonian of a region $B$ equals the bulk modular Hamiltonian of its entanglement wedge plus an area operator: $K_B=\hat A/4G_N+K_{\rm bulk}$. In the language of \cref{ch:modham}: relative entropies of boundary states for $B$ equal relative entropies of bulk states in the entanglement wedge. Almheiri, Dong and Harlow identified this structure as a quantum error-correcting code, and Harlow gave the operator-algebra version: a bulk operator in the wedge can be reconstructed on $B$ iff it commutes with the complement's algebra. Algebraically, it says: the boundary algebra $\mathcal M_B$ of $B$, restricted to the code subspace, equals the bulk algebra $\A(\mathcal E_B)$ of the entanglement wedge $\mathcal E_B$, and $\mathcal M_B'\leftrightarrow\A(\mathcal E_B')$.

\subsection{Quantum extremal surfaces and islands}

The generalized entropy is extremized by the quantum extremal surface (Engelhardt--Wall \cite{EngelhardtWall2015}): $\partial\Sgen=0$. In the crossed-product framework \cite{CPW2022,JSS2023} the stationarity of $\Sgen$ becomes the stationarity of a von Neumann entropy of a type II algebra, with the dressing providing the area operator. In evaporating black holes, the entanglement wedge of the radiation includes an \emph{island}\index{island} in the interior (Penington \cite{Penington2020}; Almheiri--Engelhardt--Marolf--Maxfield \cite{AEMM2019}): in algebraic terms, the radiation algebra, in the code subspace, contains interior operators, and the Page-curve entropy is the type II entropy of the radiation algebra with this larger wedge.

\begin{whatwrong}
The error-correction statements hold in the code subspace of semiclassical states with fixed background geometry; they cannot be extended off the code subspace, and the algebra equality is up to $1/N$ corrections. The role of type III$_1$ versus type I at finite $N$ is delicate: both appear at different levels of approximation.
\end{whatwrong}

\subsection*{Exercises}
\begin{exercise}\label{ex:ew_qec:1}
Show that if the boundary and bulk modular Hamiltonians are related by JLMS, then relative entropies on $\mathcal M_B$ and on $\A(\mathcal E_B)$ coincide on code states.
\end{exercise}
\begin{exercise}\label{ex:ew_qec:2}
Explain why a bulk operator in the entanglement wedge of $B$ commutes with all boundary operators in $B'$: consider the complementary wedge algebra and locality of the bulk.
\end{exercise}

\section{Curved spacetime and Hawking radiation}\label{ch:curved_hawking}
\skippable

\begin{physmotiv}
On a curved background with no preferred vacuum, the algebraic approach shines: one defines the algebra locally, from the equations of motion, and chooses states by a local criterion (Hadamard). The Unruh and Hawking effects are KMS statements about the horizon. The same relative entropy monotonicity that gave the GSL in \cref{ch:gen_entropy} yields the null energy-type inequalities.
\end{physmotiv}

\subsection{Algebraic QFT on curved spacetime}

On a globally hyperbolic spacetime $(\mathcal M,g)$ one assigns to each region the Weyl algebra of the classical solution space of the wave equation (\cref{ch:weyl}), with the symplectic form given by the causal propagator; locality and isotony hold by construction and there is no preferred vacuum or Hilbert space. A state is a positive functional; the physically admissible (\emph{Hadamard}) states have two-point functions with the short-distance singularity of Minkowski space; they exist on any globally hyperbolic background. This yields a notion of local quantum field theory without Poincar\'e symmetry.

\subsection{Hawking radiation as KMS}

For a Killing horizon (Rindler, Schwarzschild, de Sitter), a Hadamard state regular on the horizon (the Hartle--Hawking--Israel, Unruh or Bunch--Davies state) is KMS with respect to the horizon Killing flow at inverse temperature $\beta=2\pi/\kappa$, with $\kappa$ the surface gravity (Kay--Wald \cite{KayWald1991}). The Unruh state of a collapsing star is thermal at infinity at $T_H=\kappa/2\pi$, as the modular theory suggests.

\subsection{Relative entropy, ANEC and QNEC}

Monotonicity of relative entropy under restriction to subalgebras, applied to null-cut algebras, gives energy inequalities: the averaged null energy condition and the quantum null energy condition, which bounds $\langle T_{kk}\rangle$ below by $\hbar/2\pi$ times the second variation of the entropy with respect to the cut; these were proved in various settings (Bousso, Fisher, Leichenauer, Wall and collaborators for the conjecture and free fields; Balakrishnan, Faulkner, Khandker and Wang \cite{BFKW2019}, and Ceyhan and Faulkner \cite{CeyhanFaulkner2020}, for general proofs using relative entropy and the Connes cocycle).  In the gravitational context the same inequalities underlie the GSL (\cref{ch:gen_entropy}).

\subsection*{Exercises}
\begin{exercise}\label{ex:curved_hawking:1}
Compute the Hawking temperature of Schwarzschild from $\kappa=1/4GM$ and relate to the modular flow of the Hartle--Hawking state.
\end{exercise}
\begin{exercise}\label{ex:curved_hawking:2}
Show that for a free scalar, $S_{\rm rel}\ge0$ for null-cut algebras implies $\int\dd\lambda\,T_{\lambda\lambda}\ge0$ in a semiclassical regime (the ANEC), at the level of the formal argument in \cref{ch:gen_entropy}.
\end{exercise}


\section{Black holes and de Sitter: the sign of the thermodynamics}\label{ch:bhdS}
\begin{physmotiv}
The black-hole and de Sitter crossed products have the same form but opposite types, II$_\infty$ versus II$_1$. The difference is one sign in the thermodynamics: energy added to a black hole raises its entropy, energy added to a de Sitter static patch lowers it. This section does the arithmetic.
\end{physmotiv}

\subsection{Schwarzschild}
With $\hbar=c=k_B=1$ the Schwarzschild black hole of mass $M$ has horizon radius $r_s=2GM$, Hawking temperature and entropy
\begin{equation}
T_H=\frac1{8\pi GM},\qquad S_{\rm BH}=\frac{A}{4G}=4\pi GM^2,\qquad\frac{\dd S_{\rm BH}}{\dd M}=8\pi GM=\frac1{T_H}.
\end{equation}
The heat capacity $C=\dd M/\dd T_H=-8\pi GM^2$ is \emph{negative}: the black hole gets hotter as it loses energy. In the microcanonical sum $\int\dd M\,\ee^{S(M)}$ the entropy $S(M)$ increases without bound.

\subsection{De Sitter static patch}
With an observer of energy $E$ at the centre the horizon area is $A_{\rm dS}-\beta_{\rm dS}\cdot4G\,E$ to first order (\cref{ch:anatomy}), so
\begin{equation}
S_{\rm patch}(E)=S_{\rm dS}-\beta_{\rm dS}E,\qquad\frac{\dd S}{\dd E}=-\beta_{\rm dS}<0,
\end{equation}
with a \emph{fixed} temperature $\beta_{\rm dS}^{-1}$ (the area, and hence the entropy, \emph{decreases} when energy is placed inside the patch).

\subsection{Why the types differ}
In both cases the weight on the clock is $\ee^{S}$ times a Boltzmann factor, which in the crossed-product trace is $\int\dd x\,\ee^{x}$. For the black hole the energy variable is unbounded above and $\ee^{S(M)}$ keeps growing: the total weight is infinite, so the algebra is II$_\infty$ and there is no maximum-entropy state. For de Sitter, $E\ge0$ and the weight $\ee^{-\beta E}$ decays, so the total weight is finite (\eqref{eq:trP}) and the algebra is II$_1$.

\subsection*{Exercises}
\begin{exercise}\label{ex:bhdS:1}
Verify $T_H$, $S_{\rm BH}$ and $C$ above from $r_s=2GM$ and $T_H=\kappa/2\pi$, $\kappa=1/4GM$.
\end{exercise}
\begin{exercise}\label{ex:bhdS:2}
Show that the canonical partition function $Z(\beta)=\int\dd M\,\ee^{S(M)-\beta M}$ for a Schwarzschild black hole has its saddle at $M=\beta/8\pi G$ and that this saddle is a minimum, not a maximum, of the exponent.
\end{exercise}
\begin{exercise}\label{ex:bhdS:3}
Using $\dd M/\dd t=-1/(15360\pi G^2M^2)$ (the Stefan--Boltzmann rate for a Schwarzschild black hole radiating massless modes), find the evaporation time. Is the temperature of the final stage consistent with the semiclassical approximation?
\end{exercise}
\begin{exercise}\label{ex:bhdS:4}
For the de Sitter patch with $E\ge0$ compute the canonical partition function $\int_0^\infty\dd E\,\ee^{S_{\rm patch}(E)-\beta'E}$ for a heat bath at an arbitrary temperature $\beta'^{-1}$ and show it converges iff $\beta'>-\beta_{\rm dS}$, so that for positive temperatures it always converges.
\end{exercise}
